% R. Brown's translations of two letters to L. Breen, with his note and comments (cote 134-2): archive cote 134-2, site page 133 (PDF page 134).
% translation of the letter of 17-19 July 1975. Draft compared: ps1.tex, lines 1539-1570.
% AI-transcribed by Claude Opus 5.5 (claude-opus-5-5, adaptive thinking, effort high), Message Batches API, batch brown_full (msgbatch_01Md3i5yK1PxkpJn5GHfEExE), 2026-11-13; not yet checked by eye.
\scanpage{134-2/133}
\mspage{- 6 -}

One says that an n-category $E_n$ is ``special'' ( or \uline{n-galois}) if
it is n-equivalent to the category of locally constant (n - 1)-stacks
on a convenient topological space (or a topos), or what also turns
out, it is n-equivalent to the category of n-functors
$G_n \to (\text{n - Cat})$ , where $G_n$ is an n-groupoid. If $X$ , $G_n$ , $E_n$
correspond in this way, one calls $G_n$ the fundamental n-groupoid
of $X$ , or of $E_n$ , or that $E_n$ is the category of local n-systems
on $X$ , or on $G_n$ , or that $X$ is the geometric realisation of
$G_n$ or of $E_n$ , In analogy with the familiar case $n = 1$, $G_n$
can be interpreted as the full sub-n-category of $\underline{\mathrm{Hom}}_n(E_n,(\text{n - cat}))$
formed by the n-functors $E_n \to (\text{n-cat})$ satisfying certain
exactness properties (or what is equivalent: which commute with
finite $\varprojlim$ and arbitrary $\varinjlim$ ); but this raises the disquieting
vision of n-limits in n-categories. (N.B. The case $n = 2$ begins
to become familiar to us ...). It is prudent in all of this to
suppose that $X$ is ``locally homotopically trivial'', which ensures
the pro-simplicial set which Artin-Mazur associate to it (with
the help of nerves of hyper coverings) is essentially constant in
the ordinary homotopy category - then $X$ defines the usual
homotopy type. This is surely \uline{not} the case for the étale topos
of a scheme. In such case, the fundamental n-groupoid should be
conceived as a \uuline{pro}-\uline{n}-\uline{groupoid} (nothing surprising in that, in view
of the familiar theory of $\pi_1$), and $E_n$ as an (ind)-n-category
(the ind-structure corresponds to the exigencies of local triviality
for a variable n-stack, relative to coverings more and more fine on
X).

\margin{134}
